\begin{proposition}[Fixed support point from a strict prior average]\label{prop:step-014-fixed-extraction}
Under Proposition~\ref{prop:step-013-product-lower-bound}, fix its
same \(k,N,M,n,A\), finite prior \(\mu_{N,M}\), and ideal experiment. Then
there exists a deterministic vector
\[
\boldsymbol z^*\in\operatorname{supp}(\mu_{N,M}^{\otimes k})
\subseteq\mathcal I_N^k
\tag{8}
\]
such that
\[
\mathbb E_{\substack{S\sim(P_{\boldsymbol Q^*}^{c_{\boldsymbol t^*}})^n\\
\rho_A}}
R_{P_{\boldsymbol Q^*}}\bigl(A_{\rho_A}(S),c_{\boldsymbol t^*}\bigr)
>2^{-9}.
\tag{9}
\]
Once \(\boldsymbol z^*\) is selected, the only randomness in (9) is the
fixed-instance sample and the learner coins.
\end{proposition}

\begin{proof}
Let
\[
\mathcal S=\operatorname{supp}(\mu_{N,M}^{\otimes k}).
\tag{10}
\]
Proposition~\ref{prop:step-013-product-lower-bound} supplies a finite
probability law, so
\(\mathcal S\) is finite. For each \(\boldsymbol z\in\mathcal S\), put
\[
w_{\boldsymbol z}=\mu_{N,M}^{\otimes k}(\{\boldsymbol z\}).
\tag{11}
\]
Then \(w_{\boldsymbol z}>0\) and
\(\sum_{\boldsymbol z\in\mathcal S}w_{\boldsymbol z}=1\).

Conditional on \(\boldsymbol\Xi=\boldsymbol z\), the sample and output
experiment in Proposition~\ref{prop:step-013-product-lower-bound}
is exactly the fixed-instance experiment defining \(\Phi_A(\boldsymbol z)\).
Finite expansion of the outer prior expectation therefore gives
\[
\begin{aligned}
&\mathbb E_{\substack{
  \boldsymbol\Xi\sim\mu_{N,M}^{\otimes k}\\
  S^{\mathrm{id}}\sim(P_{\boldsymbol Q}^{c_{\boldsymbol T}})^n\\
  H^{\mathrm{id}}\sim A(S^{\mathrm{id}})}}
  R_{P_{\boldsymbol Q}}(H^{\mathrm{id}},c_{\boldsymbol T})\\
&\qquad=\sum_{\boldsymbol z\in\mathcal S}
  w_{\boldsymbol z}\Phi_A(\boldsymbol z).
\end{aligned}
\tag{12}
\]
This is finite conditional expectation. It does not assert that the
unconditional prior mixture is iid from one deterministic distribution.

The left side of (12) is strictly greater than \(2^{-9}\) by Proposition~\ref{prop:step-013-product-lower-bound}; hence
\[
\sum_{\boldsymbol z\in\mathcal S}w_{\boldsymbol z}\Phi_A(\boldsymbol z)
>2^{-9}.
\tag{13}
\]
If every \(\boldsymbol z\in\mathcal S\) had
\(\Phi_A(\boldsymbol z)\le2^{-9}\), then
\[
\sum_{\boldsymbol z\in\mathcal S}w_{\boldsymbol z}\Phi_A(\boldsymbol z)
\le\sum_{\boldsymbol z\in\mathcal S}w_{\boldsymbol z}2^{-9}
=2^{-9},
\tag{14}
\]
contradicting (13). Thus some \(\boldsymbol z^*\in\mathcal S\) has
\(\Phi_A(\boldsymbol z^*)>2^{-9}\), which is (9). \end{proof}

\begin{proposition}[Same-instance PAC contradiction]\label{prop:step-014-fixed-contradiction}
Under Proposition~\ref{prop:step-014-fixed-extraction} and Proposition~\ref{prop:step-006-pointwise-pac-expectation}, fix the
deterministic vector \(\boldsymbol z^*\) extracted in (8). Then
\[
2^{-9}<\Phi_A(\boldsymbol z^*)\le2^{-12},
\tag{15}
\]
which is impossible.
\end{proposition}

\begin{proof}
The extraction proposition gives
\[
\Phi_A(\boldsymbol z^*)>2^{-9}.
\tag{16}
\]
This lower bound concerns the fixed \(P_{\boldsymbol Q^*}\), fixed target
\(c_{\boldsymbol t^*}\), exact sample size \(n\), fixed randomized learner
\(A\), its internal coin law, and the setting-defined population 0-1 risk.
No component is changed by selecting \(\boldsymbol z^*\).

Because \(\boldsymbol z^*\in\mathcal I_N^k\), Proposition~\ref{prop:step-006-pointwise-pac-expectation}, which is the
pointwise consequence of Assumption~\ref{assump:distribution-free-realizable-pac},
applies to this same vector:
\[
\begin{aligned}
\Phi_A(\boldsymbol z^*)
&=\mathbb E_{\substack{
  S\sim(P_{\boldsymbol Q^*}^{c_{\boldsymbol t^*}})^n\\ \rho_A}}
  R_{P_{\boldsymbol Q^*}}\bigl(A_{\rho_A}(S),c_{\boldsymbol t^*}\bigr)\\
&\le\alpha_0+\beta_0
=2^{-13}+2^{-13}=2^{-12}.
\end{aligned}
\tag{17}
\]
Equations (16) and (17) concern literally the same deterministic number;
neither includes a prior average. Finally,
\[
2^{-9}=8\cdot2^{-12}>2^{-12},
\tag{18}
\]
so (16)--(17) contradict one another. \end{proof}
